\documentclass[11pt,a4paper]{article}
\usepackage[T1]{fontenc}
\usepackage[utf8]{inputenc}
\usepackage{lmodern}
\usepackage[margin=2.15cm]{geometry}
\usepackage{setspace}
\setstretch{1.08}
\setlength{\parskip}{0.25em}
\setlength{\parindent}{1em}
\usepackage{amsmath,amssymb}
\usepackage{booktabs}
\usepackage[hidelinks]{hyperref}

\title{A Physics-Informed Oxygen-Transfer Sensor\\for Brewery Aeration Fans}
\author{Short paper on innovation 2 of the thesis}
\date{}

\begin{document}
\maketitle
\thispagestyle{empty}

\begin{abstract}
\noindent
Detergents from brewery clean-in-place lower the oxygen-transfer factor \(\alpha\) that aeration controllers usually freeze. A dissolved-oxygen probe cannot separate that drop from a change in oxygen uptake: both make the reading fall, and a fan loop adds air for both. This paper specifies a soft sensor whose outputs are \(\alpha\), the transfer coefficient, and the uptake, trained with the oxygen balance and with the blower power map, not only with the probe. A health bit returns the fan law to a recipe table when the residual is bad. On a representative detergent window, using the lagged true factor instead of a quiet-week value of 0.55 raises the dissolved-oxygen minimum from 0.90 to 1.63\,mg/L. An open-loop step with uptake held fixed falls to about 0.53\,mg/L in 0.8\,h if the wrong factor is used. The sensor is justified by that floor, not by a further large cut in kilowatt-hours once the set-point is already reasonable.
\end{abstract}

\section{Introduction}

Fine-bubble aeration is sensitive to surfactants. The process-water transfer rate over the clean-water rate, the factor \(\alpha\), falls when detergents occupy the bubble interface, and the fall is sharper for fine bubbles than for coarse ones \cite{rosso2006,mueller2002}. Brewery clean-in-place liquors are a surfactant source by design \cite{simate2011,goode2013,fillaudeau2006}. Municipal soft sensors and recent physics-informed dissolved-oxygen models forecast concentration \cite{haimi2013,koksal2024,raissi2019}. They do not emit \(\alpha\) into a fan command. Off-gas hoods measure transfer directly and remain the reference \cite{redmon1983,leu2009}. They are a poor assumption as the only signal at the moment a recipe arrives.

This paper defines the sensor as useful only if a fan law can consume it. The companion paper writes that law. Here the law is the customer: each minute it inverts an oxygen-transfer curve that is proportional to \(\alpha\). A forecast of dissolved oxygen that hides a wrong \(\alpha\) inside a wrong uptake authorizes the wrong airflow.

\section{What one probe does not identify}

Ignore advection. A steady oxygen balance is
\begin{equation}
\label{eq:line}
K_La\,(\beta C^{*}-C)-\mathrm{OUR}\approx 0.
\end{equation}
Any pair \((K_La,\mathrm{OUR})\) on that line matches the same \(C\). A single airflow and a single probe cannot tell a detergent hour from a load hour. The fan curve breaks the tie if the clean-water map is known. Delivered airflow \(q\) and electrical power fix \(q\) through the blower map. The map used with this plant is polytropic compression against
\begin{equation}
\Delta p = 52 + 0.55\,(q/1000)^{2}\quad[\mathrm{kPa}],
\end{equation}
which is mostly tank submergence. The clean-water curve then fixes \(K_La\) up to \(\alpha\),
\begin{equation}
\label{eq:kla}
K_La=\alpha\,k_{\mathrm{ref}}(q/q_{\mathrm{ref}})^{0.82}\,\theta^{\,T-20},
\end{equation}
with \(k_{\mathrm{ref}}=6.5\,\mathrm{h^{-1}}\) at \(2500\,\mathrm{Nm^{3}/h}\) and \(\theta=1.024\). Once \(K_La\) and \(C\) are known, \eqref{eq:line} solves for uptake. The partial derivatives at fixed \(q\) and \(T\),
\begin{equation}
\frac{\partial K_La}{\partial\alpha}
=
k_{\mathrm{ref}}\left(\frac{q}{q_{\mathrm{ref}}}\right)^{0.82}\theta^{\,T-20}>0,
\qquad
\frac{\partial\,\mathrm{OUR}}{\partial K_La}
=
\beta C^{*}-C,
\end{equation}
are nonzero whenever the tank is not saturated. A local change in \(\alpha\) cannot hide inside the uptake. What remains is a flat direction when \(q\) and the detergent flag never move. The estimator must then hold \(\alpha\), which is a standard excitation requirement \cite{ljung1999}, not a reason to drop the balance.

At a detergent point \(K_La=6.50\,\mathrm{h^{-1}}\), \(\alpha=0.40\), and driving force \(4.60\,\mathrm{mg/L}\), the sensitivity of transfer rate to \(\alpha\) is \((6.50/0.40)\times 4.60=74.8\,\mathrm{mg\,L^{-1}\,h^{-1}}\) per unit \(\alpha\). An error of \(0.15\), the gap between a quiet-week \(0.55\) and a detergent \(0.40\), is \(11.2\,\mathrm{mg\,L^{-1}\,h^{-1}}\), more than a third of a \(30\,\mathrm{mg\,L^{-1}\,h^{-1}}\) uptake. That is the budget a sensor is asked to return.

\section{Sensor}

The deployed map reads a thirty-minute window at one-minute samples: airflow, power, header pressure, temperature, influent flow, dissolved oxygen, and the detergent flag. Thirty minutes covers several blower lags and does not cover a whole clean-in-place, so the network cannot memorize a regime instead of using the flag. Two hidden layers of 32 units feed four heads. The \(\alpha\) head is squashed into \((0.15,0.90)\). The transfer and uptake heads are positive. A fourth head predicts dissolved oxygen one step ahead and is not given to the controller. Publishing a second concentration next to a real probe invites an unaudited average.

The loss has the usual shape \cite{raissi2019,karniadakis2021,wang2021}. On a window,
\begin{equation}
\label{eq:loss}
\begin{aligned}
\mathcal{L}
&=
\frac{1}{N}\sum(\hat C-C)^{2}
+
\lambda\sum\Big(
\dot C-\widehat{K_La}(\beta C^{*}-C)+\widehat{\mathrm{OUR}}
\Big)^{2}\\
&\quad
+
\lambda\sum\big(P-P_{\mathrm{map}}(q)\big)^{2}
+
\lambda\sum\big(\widehat{K_La}-K_La^{\mathrm{curve}}(\hat\alpha,q,T)\big)^{2}.
\end{aligned}
\end{equation}
The derivative \(\dot C\) is a smoothed probe slope, not another network. The power residual stops a claimed airflow from disagreeing with the meter. The curve residual is the term that removes \eqref{eq:line}'s free direction. Training raises \(\lambda\) only after a data-only stage has fitted the probe, because the two gradients otherwise conflict \cite{wang2021,koksal2024}. The bound term penalizes \(\alpha\) moving faster than \(0.5\) per hour unless the detergent flag has changed. Diffuser wear is slow. A recipe is not.

The health bit is cleared if \(\alpha\) leaves \([0.15,0.90]\), if \(\mathcal{L}\) on the physics terms exceeds a threshold set from a commissioning week that includes one detergent day, or if the uptake head moves faster than the kinetic ceiling. A cleared bit does not stop the fans. The controller loads \(\alpha\) from the recipe table. For fermentation clean-in-place that table value is \(0.40\), already closer to the truth than \(0.55\). A stuck probe freezes airflow rather than integrating a bias against a dead reading \cite{rieger2003}. A lost calendar keeps the sensor and drops the preview. Those degradations are the interface.

Normalization is frozen at commissioning: airflow over 2500\,Nm\(^{3}\)/h, power over 80\,kW, header pressure relative to 52\,kPa, temperature relative to \(30^{\circ}\mathrm{C}\). An online mean would shift when the detergent arrives and would hide the event. Inference is a small forward pass. If it is late, the previous \(\alpha\) is kept. A fan command does not wait on a network timeout.

In the estimate reported below, the ``sensor'' is not one training seed. It is the true regime \(\alpha\), filtered with a 0.30\,h lag and a two-percent caution. That is an upper bound on what a good sensor is worth. A network whose \(\alpha\) error exceeds that filtered truth will do worse. Commissioning still needs an off-gas or portable probe on a detergent day \cite{leu2009,asce2007}. The threshold that never fires, and the threshold fit only on a quiet week, are both wrong. The first removes the fallback. The second declares every detergent day a fault and ignores the sensor exactly when it is needed.

\section{What the factor changes}

Two airflows are frozen after a detergent step, with uptake held at \(30\,\mathrm{mg\,L^{-1}\,h^{-1}}\), \(T=33^{\circ}\mathrm{C}\), and \(\beta C^{*}=6.80\). The short airflow, \(3570\,\mathrm{Nm^{3}/h}\), is the inversion of \eqref{eq:kla} at \(\alpha=0.55\). The corrected airflow, \(5250\,\mathrm{Nm^{3}/h}\), uses \(\alpha=0.40\). Dissolved oxygen starts at 1.80\,mg/L. The semi-implicit step of the oxygen balance at 0.1\,h, with uptake not allowed to fall, gives Table~\ref{tab:ol}.

\begin{table}[ht]
\centering
\caption{Open-loop dissolved oxygen (mg/L) after detergent arrives. Uptake is held at 30\,mg\,L\(^{-1}\)\,h\(^{-1}\).}
\label{tab:ol}
\small
\begin{tabular}{@{}ccc@{}}
\toprule
Time (h) & \(q=3570\) (\(\alpha\) believed 0.55) & \(q=5250\) (true \(\alpha=0.40\)) \\
\midrule
0.0 & 1.80 & 1.80 \\
0.2 & 1.08 & 2.05 \\
0.4 & 0.75 & 2.14 \\
0.8 & 0.53 & 2.18 \\
\bottomrule
\end{tabular}
\end{table}

By 0.3\,h the wrong-factor column is already under 1\,mg/L. The corrected column settles near the 2.2\,mg/L set-point. Feedback in the full week does not reach 0.53\,mg/L, because uptake falls with dissolved oxygen and because an integral bias adds air. The week minimum with \(\alpha\) locked at 0.55 is 0.90\,mg/L. With the lagged true factor it is 1.63\,mg/L. The operational floor used in the study is 1.5\,mg/L. The sensor is the difference between those two minima. Blower energy of the two predictive laws over the whole week is 9620\,kWh against 9656\,kWh. Inside the 18\,h fermentation window the live factor costs about 255\,kWh relative to the locked factor (2441 against 2186\,kWh) and buys the floor back. A report that ranks the locked factor first on kilowatt-hours in that window is ranking a loop that crossed the floor.

A proportional-integral loop at a fixed 2.5\,mg/L also stays above the floor in that window (minimum 1.87\,mg/L) and spends the difference in the quiet hours. The sensor does not replace that loop. It tells a set-point-aware inverter which \(\alpha\) to use so the set-point is reachable when the recipe, not the probe, is the first sign of detergent \cite{amand2013,rieger2014}.

\section{Commissioning check}

The threshold \(\epsilon\) and the recipe-table value of \(\alpha\) are identified on a day when the proportional-integral fallback is still in charge and the sensor is only logging. The day has to include a real detergent recipe. A threshold fit on a quiet week marks the detergent day as a fault, and the fan law then ignores \(\alpha\) at the only hour Table~\ref{tab:ol} says it matters. A threshold so loose that the health bit never clears deletes the fallback to the recipe table. The logged residual on that detergent day, compared with a portable probe or an off-gas hood, is the quantile \cite{leu2009,redmon1983}.

Two other checks stop a good dissolved-oxygen forecast from being mistaken for a good \(\alpha\). First, the curve residual in \eqref{eq:loss} is inspected on the same day. If the emitted \(K_La\) and the curve at the emitted \(\alpha\) disagree by more than about fifteen percent, the health bit is treated as clear even if the network set it. Second, the fan inversion is evaluated twice on the logged airflow: once with the emitted \(\alpha\) and once with 0.55. The gap should be of the order of the 3570 against 5250\,Nm\(^{3}\)/h in Table~\ref{tab:ol} when the recipe is the fermentation clean. If the two inversions agree, the sensor has not moved \(\alpha\), and the network is not yet doing the job this paper defined. Agreement of the one-step dissolved-oxygen head with the probe is not a substitute for that gap. Equation~\eqref{eq:line} can be fitted by moving uptake alone.

The sensor is installed last. Constant-speed headers move to a dissolved-oxygen loop first. The recipe clock and the set-point table come second, with \(\alpha\) from the table. The network closes the inversion only after the detergent-day log exists. That order is forced by the numbers: the large energy cut does not need the network, and the dissolved-oxygen floor during detergent does not survive a wrong \(\alpha\).

\section{Conclusion}

The oxygen-transfer factor during brewery clean-in-place is a measured disturbance, not a calibrated constant. A small residual network can emit it from fan power, airflow, and a dissolved-oxygen probe if the clean-water curve is enforced, so that uptake cannot trade off against \(\alpha\). On the representative window, replacing a quiet-week factor of 0.55 with the lagged true factor of about 0.40 lifts the dissolved-oxygen minimum from 0.90 to 1.63\,mg/L at essentially the same weekly blower energy. The open-loop column shows why the weekly minimum is the mitigated error, not the raw one. The result is an integration with a filtered true factor standing in for a commissioned sensor. It is not a field residual, and it is not a claim that the sensor saves a further large fraction of fan electricity once the airflow set-point is already right.

\begin{thebibliography}{99}
\bibitem{rosso2006} D.~Rosso and M.~K. Stenstrom. Surfactant effects on \(\alpha\)-factors in aeration systems. \textit{Water Research}, 40:1397--1404, 2006.
\bibitem{mueller2002} J.~A. Mueller, W.~C. Boyle, and H.~J. P{\"o}pel. \textit{Aeration: Principles and Practice}. CRC Press, 2002.
\bibitem{simate2011} G.~S. Simate et al. The treatment of brewery wastewater for reuse: State of the art. \textit{Desalination}, 273:235--247, 2011.
\bibitem{goode2013} K.~R. Goode, K.~Asteriadou, P.~T. Robbins, and P.~J. Fryer. Fouling and cleaning studies in the food and beverage industry classified by cleaning type. \textit{Comprehensive Reviews in Food Science and Food Safety}, 12:121--143, 2013.
\bibitem{fillaudeau2006} L.~Fillaudeau, P.~Blanpain-Avet, and G.~Daufin. Water, wastewater and waste management in brewing industries. \textit{J. Cleaner Production}, 14:463--471, 2006.
\bibitem{haimi2013} H.~Haimi, M.~Mulas, F.~Corona, and R.~Vahala. Data-derived soft-sensors for biological wastewater treatment plants: An overview. \textit{Environmental Modelling and Software}, 47:88--107, 2013.
\bibitem{koksal2024} E.~S. Koksal and E.~Aydin. A hybrid approach of transfer learning and physics-informed modelling for dissolved oxygen in an industrial wastewater treatment plant. \textit{Chemical Engineering Science}, 121088, 2024.
\bibitem{raissi2019} M.~Raissi, P.~Perdikaris, and G.~E. Karniadakis. Physics-informed neural networks. \textit{Journal of Computational Physics}, 378:686--707, 2019.
\bibitem{redmon1983} D.~Redmon, W.~C. Boyle, and L.~Ewing. Oxygen transfer efficiency measurements in mixed liquor using off-gas techniques. \textit{J. Water Pollution Control Federation}, 55:1338--1347, 1983.
\bibitem{leu2009} S.-Y. Leu, D.~Rosso, L.~E. Larson, and M.~K. Stenstrom. Real-time aeration efficiency monitoring in the activated sludge process. \textit{Water Environment Research}, 81:2471--2481, 2009.
\bibitem{ljung1999} L.~Ljung. \textit{System Identification: Theory for the User}. Prentice Hall, 2nd ed., 1999.
\bibitem{karniadakis2021} G.~E. Karniadakis et al. Physics-informed machine learning. \textit{Nature Reviews Physics}, 3:422--440, 2021.
\bibitem{wang2021} S.~Wang, Y.~Teng, and P.~Perdikaris. Understanding and mitigating gradient pathologies in physics-informed neural networks. \textit{SIAM J. Scientific Computing}, 43:A3055--A3081, 2021.
\bibitem{rieger2003} L.~Rieger, J.~Alex, S.~Winkler, M.~Boehler, M.~Thomann, and H.~Siegrist. Progress in sensor technology --- progress in process control? \textit{Water Science and Technology}, 47(2):103--112, 2003.
\bibitem{asce2007} ASCE. \textit{Measurement of Oxygen Transfer in Clean Water}, ASCE/EWRI 2-06. Reston, 2007.
\bibitem{amand2013} L.~{\AA}mand, G.~Olsson, and B.~Carlsson. Aeration control --- a review. \textit{Water Science and Technology}, 67:2374--2397, 2013.
\bibitem{rieger2014} L.~Rieger, R.~M. Jones, P.~L. Dold, and C.~B. Bott. Ammonia-based feedforward and feedback aeration control in activated sludge processes. \textit{Water Environment Research}, 86:63--71, 2014.
\end{thebibliography}

\end{document}
